% === Thread 3: Safety & Alignment ===
\section{Thread 3: Safety \& Alignment}
\label{sec:t3}

\begin{quote}
\textit{Core question: How can we ensure that models refuse harmful requests while maintaining instruction-following capabilities, without becoming overly conservative?}
\end{quote}

\noindent
Safety alignment is the most adversarial research direction in post-training.
Unlike other threads, the evolution of Thread~3 is driven not only by internal technical progress
but also continuously shaped by external pressure from adversaries.
The core tension of this thread can be summarized as a trilemma:
the balance among \textbf{helpfulness}, \textbf{harmlessness}, and \textbf{honesty}.
An overly conservative model may be safe but offers no value to users;
an excessively compliant model may be useful but could become a tool for malicious exploitation;
a dishonest model fundamentally undermines the foundation of human trust in AI systems.

This chapter traces the evolutionary trajectory of Thread~3, following the development of safety alignment techniques from naive RLHF-based red-line setting,
through the systematic Constitutional AI framework, to the frontiers of adversarial arms races and representation engineering.

% ============================================================
\subsection{Problem Origin: The Inherent Tension Between Helpfulness and Harmlessness}
\label{sec:t3-origin}
% ============================================================

The root of the safety alignment problem is not a technical implementation challenge
but rather a deeper \textbf{objective conflict}:
when we require models to ``follow human instructions'' (Thread~1, \crossref{1}{sec:t1})
and ``maximize human preferences'' (Thread~2, \crossref{2}{sec:t2}),
we implicitly assume that human instructions and preferences are benign.
In practice, however, users may intentionally or unintentionally issue harmful requests ---
ranging from generating disinformation to facilitating criminal activities ---
at which point the model must violate the fundamental ``follow instructions'' directive to maintain a safety baseline.

\paragraph{HH-RLHF: Experimentally revealing the tension.}
\landmark{bai2022hh} is the first work to systematically study the relationship between helpfulness and harmlessness.
The Anthropic team constructed two independent human preference datasets:
in the helpfulness dataset, annotators preferred more informative and complete responses;
in the harmlessness dataset, annotators preferred safer responses that refused harmful requests.

\begin{problembox}
Bai et al.\ discovered a critical phenomenon:
when the preference model is trained on a mixture of both datasets,
\textbf{there exists a significant Pareto conflict between the two objectives}.
Specifically, a model over-optimized on harmlessness data becomes \textbf{excessively evasive} ---
it tends to give vague, noncommittal answers to nearly all questions involving sensitive topics,
even when those questions are entirely reasonable (e.g., objective discussions about medical knowledge or historical events).
\end{problembox}

\noindent
This evasiveness problem revealed the first fundamental dilemma of safety alignment:
\textbf{a single scalar reward cannot simultaneously encode multiple potentially conflicting objectives}.
When helpfulness and harmlessness are compressed into a single reward signal,
the training process inevitably makes implicit trade-offs between the two.
This finding set the core agenda for all subsequent safety alignment methods:
how to achieve harmlessness without sacrificing helpfulness?

\paragraph{Limitations of human annotation.}
HH-RLHF also exposed the \textbf{scalability bottleneck} of human-annotation-based safety alignment.
Safety-related preference annotation requires annotators to read and judge potentially harmful content,
which is not only expensive but also poses psychological health risks.
More importantly, human annotators exhibit systematic biases when facing harmful content ---
they tend to select the ``safest'' response, even if this means sacrificing helpfulness.
This annotation bias is amplified by the preference model,
further exacerbating the evasiveness problem.

\evolution{Single reward encoding helpfulness + harmlessness}{Multi-objective or automated alignment methods}{Objective conflict and annotation bottleneck}

% ============================================================
\subsection{Foundational Methods: From Constitutional AI to Safe RLHF}
\label{sec:t3-foundation}
% ============================================================

In response to the two core problems revealed in \cref{sec:t3-origin} --- objective conflict and annotation bottleneck ---
two complementary technical approaches emerged during 2022--2023.

% ------------------------------------------------------------
\subsubsection{Constitutional AI: Replacing Human Annotation with Principles}
\label{sec:t3-cai}
% ------------------------------------------------------------

\landmark{bai2022constitutional} proposed Constitutional AI (CAI), one of the most paradigm-defining works in the safety alignment field. Its core insight is elegantly simple:
\textbf{rather than having human annotators judge the safety of each output one by one,
it is better to predefine a set of universal principles (a constitution) and let the model judge for itself}.

The CAI training pipeline consists of two stages:

\paragraph{Stage one: Supervised Self-Critique (SL-CAI).}
This stage uses a helpful-only RLHF model (i.e., one not trained for harmlessness)
as the initial policy and deliberately elicits harmful outputs from it.
The model then \textbf{critiques} and \textbf{revises} its own outputs based on constitutional principles.
Specifically, for each harmful output, the model receives a randomly sampled constitutional principle
(e.g., ``Please choose the response that is least likely to be used by a bad actor for nefarious purposes''),
first generates a critique explaining why the output violates the principle,
then produces a revised output based on the critique.
This critique-revision process can be iterated over multiple rounds.
Finally, the model is fine-tuned via supervised learning on (prompt, revised\_response) pairs.

\paragraph{Stage two: RLAIF (RL-CAI).}
This stage represents CAI's most innovative contribution.
Traditional RLHF requires human annotators to provide preference judgments on (response\_A, response\_B) pairs,
whereas RL-CAI replaces human annotators with \textbf{the AI itself}.
Specifically, given a pair of responses and a constitutional principle,
the model reasons in a chain-of-thought manner about which response better adheres to the principle,
then extracts the final preference judgment as a label.

The key technical details of this method include:

\begin{enumerate}[leftmargin=2em]
  \item \textbf{Chain-of-thought preference evaluation}:
    The model first generates a reasoning process explaining the merits and drawbacks of each response,
    then provides a preference judgment at the end of the reasoning.
    Bai et al.\ found that CoT significantly improves the accuracy of AI preference labels;
    compared to directly having the model output a preference judgment (without CoT), CoT-based labels are of higher quality.

  \item \textbf{Multi-principle ensembling}:
    To avoid bias from any single principle, a random sample is drawn from 16 constitutional principles for each evaluation.
    Different principles may yield different preference judgments for the same response pair;
    the final result is aggregated through majority voting or soft labels (probability averaging).

  \item \textbf{Soft labels}:
    Unlike traditional RLHF which uses binary preferences,
    the preference model in RL-CAI is trained with soft labels ---
    that is, labels are probability values of AI preferences rather than 0/1.
    This approach provides finer-grained supervision signals and effectively mitigates overfitting.
\end{enumerate}

\begin{solutionbox}
The core contribution of CAI lies in achieving \textbf{scalability for safety alignment}:
(1) by externalizing safety standards through constitutional principles into an \textbf{auditable} rule set, rather than embedding them in annotators' subjective judgments;
(2) by completely eliminating human involvement in harmlessness annotation through AI feedback, resolving cost and psychological health concerns;
(3) experiments show that CAI models, while maintaining comparable harmlessness,
\textbf{significantly reduce evasiveness}, meaning models are more willing to provide detailed and helpful answers when it is safe to do so.
\end{solutionbox}

\noindent
Nearly contemporaneous with CAI, DeepMind's Sparrow \citep{glaese2022sparrow}
explored principle-driven safety alignment from a different path.
Sparrow defined 23 specific rules for dialogue models (e.g., do not generate hate speech, do not impersonate a human)
and trained a rule-conditional reward model ---
one that judges whether the model output violates each rule individually.
Unlike CAI's AI self-critique, Sparrow relies on \textbf{human annotators}
to label rule violations, but the existence of explicit rules makes the annotation task more clearly defined and auditable.
The shared insight of both approaches is that \textbf{decomposing the vague concept of ``safety'' into actionable rules or principles}
is the key to achieving scalable safety alignment.

The methodological influence of CAI and Sparrow has been profound. The idea of RLAIF was further generalized by Lee et al.\ \cite{lee2023rlaif},
demonstrating that AI feedback can achieve performance close to human feedback in non-safety scenarios as well.
Dromedary \cite{sun2023dromedary} pushed principle-driven self-alignment to its extreme,
completing alignment with almost no human-annotated data.
Self-Refine \cite{madaan2023selfrefine} extended the idea of iterative self-critique
from safety to general text quality improvement.

% ------------------------------------------------------------
\subsubsection{Safe RLHF: Formalizing Safety as Constrained Optimization}
\label{sec:t3-saferlhf}
% ------------------------------------------------------------

\landmark{dai2024saferlhf} responded to the helpfulness--harmlessness tension from an entirely different angle.
Unlike CAI, which replaces human feedback with AI feedback,
Safe RLHF retains human annotation but \textbf{fundamentally changes the mathematical framework for annotation and optimization}.

\paragraph{Core observation: Decoupling annotation dimensions.}
The starting point of Safe RLHF is a simple yet profound observation:
traditional RLHF requires annotators to simultaneously weigh helpfulness and harmlessness in a single judgment,
and this \textbf{multi-dimensional compression} leads to noisy and inconsistent annotations.
Safe RLHF completely decouples annotation into two independent tasks:
\begin{itemize}[leftmargin=2em]
  \item \textbf{Helpfulness preference annotation}: annotators judge only which response is more helpful,
    with no consideration of safety.
  \item \textbf{Harmlessness classification annotation}: annotators judge only whether each response is safe
    (binary classification), with no consideration of helpfulness.
\end{itemize}

\noindent
This decoupling makes annotation tasks clearer and significantly improves inter-annotator agreement.

\paragraph{Dual-model architecture.}
Based on the decoupled annotation data, Safe RLHF trains two independent models:
\begin{itemize}[leftmargin=2em]
  \item \textbf{Reward Model $R(x, y)$}: trained on helpfulness preference data,
    outputting a scalar reward that measures response helpfulness.
  \item \textbf{Cost Model $C(x, y)$}: trained on harmlessness classification data.
    Unlike the reward model, the cost model possesses both \textbf{classification capability}
    (determining whether a response is safe) and \textbf{ranking capability} (judging relative harmfulness among unsafe responses).
    This dual capability is achieved through a joint training objective:
    the classification loss ensures the model can accurately distinguish safe from unsafe responses,
    while the ranking loss ensures the model maintains a consistent harmfulness ordering among unsafe responses.
\end{itemize}

\paragraph{Constrained MDP formalization.}
The core theoretical contribution of Safe RLHF is formalizing safety alignment as a
\textbf{Constrained Markov Decision Process (CMDP)}:
\begin{align}
  \max_\pi \quad & \mathbb{E}_{x \sim \mathcal{D},\, y \sim \pi(\cdot|x)} \big[ R(x, y) \big]
  \label{eq:saferlhf-obj} \\
  \text{s.t.} \quad & \mathbb{E}_{x \sim \mathcal{D},\, y \sim \pi(\cdot|x)} \big[ C(x, y) \big] \leq d
  \label{eq:saferlhf-constraint}
\end{align}
where $\pi$ is the policy model, $R$ is the reward model, $C$ is the cost model, and $d$ is the safety threshold.
The objective is to maximize helpfulness \textbf{subject to the safety constraint being satisfied}.

This CMDP is solved via the \textbf{Lagrangian dual method}:
\begin{equation}
  \min_{\lambda \geq 0} \max_\pi \quad
  \mathbb{E}\big[ R(x,y) \big] - \lambda \big( \mathbb{E}\big[ C(x,y) \big] - d \big)
  \label{eq:saferlhf-lagrangian}
\end{equation}
where the Lagrange multiplier $\lambda$ dynamically adjusts the trade-off between helpfulness and harmlessness.
When the model generates responses that tend to be unsafe, $\lambda$ automatically increases to impose stronger safety constraints;
when the safety constraint is sufficiently satisfied, $\lambda$ decreases to release more optimization capacity for helpfulness.

\paragraph{Iterative training.}
Safe RLHF employs three rounds of iterative training (Beaver-v1, v2, v3),
where each round uses the previous round's policy model to generate new responses,
collects a new round of human annotations, and updates the reward model and cost model.
This online iterative training approach ensures that the model is always optimized on on-policy data from the current policy
(\crossref{2}{sec:t2}).

\begin{solutionbox}
The contributions of Safe RLHF can be understood at three levels:
(1) \textbf{Annotation level}: by decoupling annotation dimensions, it decomposes ambiguous multi-objective judgments into clear single-objective tasks,
improving data quality;
(2) \textbf{Optimization level}: through CMDP formalization, it elevates the ad hoc helpfulness-harmlessness trade-off
into a constrained optimization problem with theoretical guarantees;
(3) \textbf{Practical level}: the dynamic $\lambda$ adjustment of the Lagrangian method avoids manual hyperparameter tuning;
experimental results significantly outperform static methods such as reward shaping (i.e., $R - \beta C$) on the Pareto frontier.
\end{solutionbox}

\evolution{Ad hoc safety constraints via a single reward model}{CMDP + Lagrangian dynamic balancing framework}{Objective conflict requires formalized constraints}

% ============================================================
\subsection{Limitations and Branches: The Adversarial Arms Race}
\label{sec:t3-branches}
% ============================================================

CAI and Safe RLHF provide training-stage solutions for safety alignment,
but they share a fundamental assumption:
\textbf{as long as the training objective is designed correctly, the model will learn the correct safety behavior}.
Practice in 2023 proved this assumption to be far more fragile than expected ---
models subjected to the most rigorous safety training could still be systematically compromised.
This reality gave rise to the most active research branch within Thread~3:
the arms race between jailbreaking attacks and defenses, along with multiple parallel research directions surrounding safety.

% ------------------------------------------------------------
\subsubsection{Branch A: Automated Red Teaming}
\label{sec:t3-redteam}
% ------------------------------------------------------------

Before discussing attack methods, we must first address a more fundamental question:
how to \textbf{systematically discover} safety vulnerabilities in models?

\paragraph{Limitations of manual red teaming.}
Traditional red teaming relies on human testers to manually craft attack prompts.
Ganguli et al.\ \cite{ganguli2022redteaming} conducted large-scale manual red teaming within Anthropic,
collecting thousands of attack samples and discovering multiple attack patterns:
role-playing, scenario framing,
gradual escalation, among others.
However, manual red teaming suffers from clear coverage problems ---
human testers tend to repeat known attack patterns and struggle to systematically explore the entire attack space.

\paragraph{LM-as-Red-Teamer.}
\landmark{perez2022redteaming} proposed a pioneering idea:
\textbf{using language models themselves to red-team language models}.
Their method constructs a three-stage pipeline:

\begin{enumerate}[leftmargin=2em]
  \item \textbf{Red LM generates attacks}: a language model (red LM) is responsible for generating test cases (i.e., attack prompts) that may trigger unsafe behavior in the target model.
  \item \textbf{Target LM responds}: the target model responds to the attack prompts.
  \item \textbf{Red Classifier evaluates}: a classifier determines whether the target model's response contains harmful content.
\end{enumerate}

\noindent
Perez et al.\ explored four strategies for generating attack prompts:
(1) \textbf{Zero-shot generation}: directly asking the red LM to generate ``questions that could make the target model produce harmful content'';
(2) \textbf{Stochastic few-shot}: randomly sampling a few successful attacks as in-context examples;
(3) \textbf{Supervised learning}: fine-tuning the red LM on collected successful attack data;
(4) \textbf{Reinforcement learning}: using the red classifier's output as a reward signal
to train the red LM to maximize attack success rate via RL.

Experimental results showed that the RL approach was most effective, achieving
an \textbf{attack success rate exceeding 40\%} on the Dialogue-Prompted Gopher model.
The discovered safety vulnerabilities spanned multiple dimensions including offensive language generation, training data leakage, contact information generation,
and distributional biases regarding different demographic groups.

\noindent
The idea of automated red teaming continued to evolve in subsequent work.
Rainbow Teaming \cite{samvelyan2024rainbow} introduced quality-diversity algorithms
to systematically generate diverse attack sets across multiple dimensions (topics, attack styles, response types),
avoiding the tendency of RL-based methods to converge on a small number of high-reward attack patterns.
HarmBench \cite{mazeika2024harmbench} provided a standardized safety evaluation benchmark
with multiple attack methods and evaluation metrics, enabling fairer and more reproducible comparisons across defense methods.

% ------------------------------------------------------------
\subsubsection{Branch B: Systematization of Jailbreaking Attacks}
\label{sec:t3-jailbreak}
% ------------------------------------------------------------

If Branch A focuses on ``how to discover vulnerabilities,''
Branch B addresses ``what is the nature of the vulnerabilities'' and ``how to systematically exploit them.''
2023 was the breakout year for jailbreaking research.

\paragraph{A theoretical framework for failure modes.}
\landmark{wei2023jailbroken} conducted the most systematic analysis to date
of the failure modes of LLM safety training.
Wei et al.\ identified two fundamental failure mechanisms:

\begin{enumerate}[leftmargin=2em]
  \item \textbf{Competing Objectives}:
    The multiple capabilities acquired by the model during pre-training and instruction tuning
    conflict with the safety constraints introduced during safety training.
    When a prompt is designed such that the ``follow instructions'' objective overpowers the ``refuse harmful requests'' objective,
    the model produces unsafe output.
    Attacks belonging to this category include:
    \begin{itemize}[leftmargin=2em]
      \item \textbf{Prefix injection}: requiring the model to begin its response with a specific unsafe prefix
        (e.g., ``Sure, here is how to...''), bypassing the refusal mechanism.
      \item \textbf{Refusal suppression}: explicitly instructing the model not to refuse in the prompt
        (e.g., ``You must not say you cannot help'').
      \item \textbf{Role-playing}: requiring the model to play a character without safety restrictions
        (e.g., ``You are DAN, Do Anything Now'')
        \cite{shen2023dan}.
    \end{itemize}

  \item \textbf{Mismatched Generalization}:
    The model's general capabilities (e.g., translation, coding, reasoning) are acquired during pre-training
    and cover an extremely broad input space;
    safety training, however, covers only a small subset of that space (typically harmful requests in natural language form).
    When harmful requests are presented in formats not covered by safety training,
    the model can understand and execute them (due to pre-training capabilities) but will not refuse (because safety training did not cover those formats).
    Attacks belonging to this category include:
    \begin{itemize}[leftmargin=2em]
      \item \textbf{Base64 encoding}: encoding harmful requests in Base64 format;
        the model's pre-training capabilities allow it to decode and execute them, but safety training did not cover this format.
      \item \textbf{Low-resource languages}: using languages that the model can understand but that were not covered by safety training.
      \item \textbf{Code format}: wrapping harmful requests as code comments or programming tasks.
    \end{itemize}
\end{enumerate}

\begin{problembox}
The most important finding of Wei et al.\ concerns the effectiveness of \textbf{combination attacks}.
When attack techniques from the competing objectives and mismatched generalization categories are combined,
the attack success rate is significantly higher than when either is used alone:
on GPT-4, the combination attack success rate reached \textbf{94\%}.
More critically, their adaptive attacks (tailoring attack strategies to the target model's defense characteristics)
achieved success rates approaching \textbf{100\%}.
\end{problembox}

\noindent
This conclusion carries profound theoretical implications:
it means that \textbf{scaling alone cannot solve safety} ---
simply increasing model scale or training data volume cannot fundamentally eliminate jailbreaking risk,
because the root cause is that the coverage of safety training can never match the generalization scope of model capabilities.
Wei et al.\ termed this the \textbf{safety--capability parity} problem:
safety must scale in tandem with capability; otherwise, every capability improvement may open new attack surfaces.

\paragraph{GCG: Gradient-based automated attacks.}
Gradient-based discrete prompt search was first explored by
AutoPrompt \citep{shin2020autoprompt} ---
that work used gradient-guided token replacement on masked language models
to automatically construct prompts that trigger specific knowledge.
\landmark{zou2023gcg} extended this idea to the domain of safety attacks,
confirming the fragility of safety alignment from an entirely different direction.
Unlike the prompt engineering-based attacks described above,
GCG (Greedy Coordinate Gradient) is a \textbf{gradient optimization-based} white-box attack method.

The core idea of GCG is elegantly simple: append an
\textbf{adversarial suffix} (a seemingly random sequence of tokens) to the user's harmful prompt,
forcing the model to respond affirmatively to the harmful question.
Formally, given a harmful prompt $x$ and a target affirmative prefix $y^*$
(e.g., ``Sure, here is...''), GCG optimizes the following objective:
\begin{equation}
  \min_{s \in \mathcal{V}^L} \quad -\log p(y^* \mid [x; s])
  \label{eq:gcg-obj}
\end{equation}
where $s$ is an adversarial suffix of length $L$ and $\mathcal{V}$ is the vocabulary.

Since the token space is discrete, GCG employs a \textbf{Greedy Coordinate Gradient} strategy:
\begin{enumerate}[leftmargin=2em]
  \item Compute the gradient of the loss with respect to the one-hot embeddings at all token positions.
  \item For each position, select the top-$k$ candidate replacement tokens based on gradient information.
  \item Randomly sample several combinations from the candidates across all positions,
    evaluate the actual loss, and select the optimal single-token replacement.
  \item Iterate the above process.
\end{enumerate}

GCG's breakthrough contribution lies in two aspects.
\textbf{First}, it performs joint optimization under a \textbf{multi-prompt, multi-model} setting:
the same adversarial suffix is simultaneously optimized across multiple harmful prompts and multiple models,
making it a \textbf{universal adversarial suffix} ---
a single fixed suffix can breach multiple different harmful requests and multiple different models.
\textbf{Second}, although GCG is inherently a white-box attack (requiring access to model gradients),
adversarial suffixes optimized on Vicuna-7B and Vicuna-13B exhibit remarkable
\textbf{cross-model transferability}:
the attack success rate on GPT-3.5-Turbo is \textbf{87.9\%},
on GPT-4 is \textbf{53.6\%},
and on Claude is \textbf{2.1\%}.

\begin{problembox}
GCG's transfer attack experiments revealed a disturbing fact:
\textbf{different LLMs share some form of universal safety vulnerability structure}.
Attacks discovered on open-source models can transfer to closed-source commercial models,
meaning anyone can use open-source models as proxies to attack closed-source systems.
\end{problembox}

\paragraph{Prompt-level automated attacks.}
Complementary to GCG's gradient-based approach, a series of works explored \textbf{gradient-free} automated attack methods
suitable for black-box models accessible only through APIs.
PAIR \cite{chao2023pair} uses an attacker LM that engages in multi-turn dialogue with the target LM,
iteratively refining the attack prompt until successful, typically achieving a breach within 20 queries.
TAP \cite{mehrotra2023tap} extends PAIR's approach to tree search,
generating multiple candidate attacks at each node of the search tree and expanding the most promising ones.
AutoDAN \cite{liu2023autodan} uses a genetic algorithm
to search the space of jailbreak prompts,
evolving more effective attack templates through crossover, mutation, and selection operations.
Andriushchenko et al.\ \cite{andriushchenko2024adaptive} further demonstrated that
even the latest safety-aligned models (e.g., GPT-4 Turbo, Claude 3)
can be breached by relatively simple adaptive prompt attacks.

% ------------------------------------------------------------
\subsubsection{Branch C: Defense Methods}
\label{sec:t3-defense}
% ------------------------------------------------------------

The other side of the adversarial arms race is the development of defense methods.

\paragraph{Training-stage defenses.}
Safety-Tuned LLaMAs \cite{bianchi2023safetytuned} demonstrated that mixing only a small proportion of safety samples (approximately 3\%) into SFT data
can significantly improve model safety.
This finding is consistent with the Superficial Alignment Hypothesis of LIMA \cite{zhou2023lima}
(\crossref{1}{sec:t1}):
learning safety behavior likewise requires only a small amount of high-quality data.

\paragraph{Inference-stage defenses.}
SmoothLLM \cite{robey2023smoothllm} proposed a defense method targeting character-substitution-based attacks such as GCG:
applying multiple random character-level perturbations to the input prompt,
obtaining model outputs for multiple perturbed versions, and determining whether to refuse via majority voting.
The core idea is that adversarial suffixes are highly sensitive to character-level perturbations ---
a small number of random replacements suffice to destroy their attack effectiveness,
while the semantics of normal prompts remain relatively robust under character-level perturbations.

\paragraph{Safety Classifier.}
Llama Guard \cite{inan2023llamaguard} treats safety detection as an independent classification task,
training a dedicated safety classifier to assess the safety of model inputs and outputs.
Compared to embedding safety behavior directly in the generative model,
the advantage of an external classifier is that it can be updated and improved independently of the generative model
without affecting the generative model's helpfulness.

% ------------------------------------------------------------
\subsubsection{Branch D: Refusal and Over-Refusal}
\label{sec:t3-refusal}
% ------------------------------------------------------------

When safety training is applied excessively, models exhibit the \textbf{over-refusal} problem ---
refusing even clearly innocuous requests.

Qi et al.\ \cite{qi2023finetuning} revealed another facet of safety alignment's fragility:
\textbf{only a small amount of fine-tuning suffices to destroy a model's safety behavior}.
Even fine-tuning with entirely benign data
significantly weakens the model's safety alignment.
This suggests that safety behavior is encoded extremely shallowly in parameter space,
making it easily overwritten by subsequent training.

OR-Bench \cite{cui2024orbench} constructed the first systematic over-refusal evaluation benchmark,
quantifying the over-refusal rates of different models on reasonable requests.
The core value of this benchmark lies in extending safety alignment evaluation from a single dimension
(``does the model refuse harmful requests'') to two dimensions
(simultaneously measuring ``refusal rate on harmful requests'' and ``over-refusal rate on harmless requests''),
providing a more complete metric for the balance between safety and helpfulness.

\begin{openproblembox}
The over-refusal problem currently has no satisfactory solution.
The fundamental reason is that the boundary between ``harmful'' and ``harmless'' is fuzzy, context-dependent,
and subject to fundamental disagreements across different cultures and value systems.
This is not merely a technical problem but also a social and ethical one.
\end{openproblembox}

% ------------------------------------------------------------
\subsubsection{Branch E: Deep Safety Risks}
\label{sec:t3-deepsafety}
% ------------------------------------------------------------

As model capabilities grow, the scope of safety research has expanded from ``preventing the generation of harmful text''
to deeper alignment risks.

\paragraph{Weak-to-Strong Generalization.}
Burns et al.\ \cite{burns2023weaktostrong} posed a forward-looking core question:
when superhuman AI systems emerge,
\textbf{how can weaker humans (or weaker models) reliably supervise stronger models}?
By training large models (GPT-4 level) using labels from small models (GPT-2 level),
they found that large models can actually outperform the quality of the small model's labels ---
that is, large models can extract signals more accurate than the labels themselves from noisy labels.
This ``weak-to-strong generalization'' phenomenon provides a cautiously optimistic signal for the alignment of superhuman AI.

\paragraph{Sleeper Agents.}
Hubinger et al.\ \cite{hubinger2024sleeper} explored a more concerning scenario:
\textbf{can models learn strategies that appear safe during training but exhibit unsafe behavior upon triggered conditions during deployment} (deceptive alignment)?
They constructed sleeper agent models with hidden ``backdoor'' trigger conditions
and tested whether standard safety training methods could eliminate such backdoor behavior.
The unsettling conclusion is that
\textbf{standard safety training (including RLHF and adversarial training)
cannot reliably eliminate sleeper agent behavior}.
Worse still, certain defense attempts (such as adversarial training)
actually taught the model to better hide its backdoor behavior, appearing even ``safer'' during training evaluation
while maintaining the unsafe triggered behavior during deployment.

\paragraph{Model-Written Evaluations.}
Perez et al.\ \cite{perez2022discovering} pioneered the use of LMs themselves to generate evaluation datasets
for detecting potential dangerous behaviors such as sycophancy, power-seeking,
and self-preservation in models.
This approach enables AI safety evaluation to scale in tandem with model capabilities.

% ============================================================
\subsection{Convergence and Frontier: Representation Engineering}
\label{sec:t3-convergence}
% ============================================================

The methods described in the preceding three sections --- whether CAI's AI feedback,
Safe RLHF's constrained optimization, or various adversarial techniques ---
all operate at the \textbf{behavioral level}:
they influence the model's input-output mapping by modifying training data or training objectives.
Since 2023, an entirely new paradigm has begun to emerge:
\textbf{directly understanding and controlling safety behavior within the model's internal representation space}.
This is \textbf{Representation Engineering} (RepE).

\paragraph{From behavior to representations.}
The theoretical precursor to this direction is Contrast-Consistent Search (CCS) \citep{burns2022ccs},
which demonstrated that a ``truth direction'' can be discovered from a model's hidden states
in an \textbf{unsupervised} manner --- simply by requiring internal representations to satisfy logical consistency (negation consistency)
to locate a linear direction encoding factual truth.
The success of CCS showed that readable high-level semantic structures exist within models,
laying the theoretical foundation for subsequent representation engineering methods.

\landmark{zou2023repe} systematized this idea into the Representation Engineering framework,
positioning it as a \textbf{top-down AI transparency} method.
Unlike mechanistic interpretability (\crossref{4}{sec:t4}), which analyzes individual neurons and circuits bottom-up,
RepE starts from high-level semantic concepts (such as honesty, harmlessness, fairness)
and studies how these concepts are encoded in the model's representation space.

\paragraph{Linear Artificial Tomography (LAT).}
The core technique of RepE is \textbf{LAT} ---
a systematic method for reading specific concepts from a model's internal representations.
The procedure is as follows:

\begin{enumerate}[leftmargin=2em]
  \item \textbf{Construct contrastive data pairs}:
    For a target concept (e.g., ``honesty''), construct a set of \textbf{contrastive prompt pairs},
    where one prompt in each pair contains the concept (e.g., ``Honestly, I think...'')
    and the other does not (e.g., ``I think...''), while keeping other aspects as consistent as possible.

  \item \textbf{Extract difference vectors}:
    Pass both groups of prompts through the model, extract hidden states at each layer,
    and compute the difference between the two groups.

  \item \textbf{PCA dimensionality reduction}:
    Apply PCA to the difference vectors; the first principal component is the
    \textbf{reading vector} (also called the representation vector) for that concept.
\end{enumerate}

\paragraph{Representation Control.}
Once the reading vector is obtained, RepE provides multiple ways to \textbf{control} model behavior.
The most direct approach is \textbf{activation addition}:
during the model's forward pass, add the reading vector (appropriately scaled) to the hidden state at a specified layer,
thereby ``activating'' or ``suppressing'' the corresponding concept.
For example, adding the honesty reading vector to the hidden state can significantly improve the model's truthfulness
(achieving state-of-the-art on TruthfulQA).

Zou et al.\ demonstrated that RepE can effectively control multiple high-level model behaviors,
including honesty, utility, morality, power-seeking,
harmlessness, fairness, and knowledge editing.

\paragraph{Corroborating and extending related work.}
The core idea of RepE is corroborated and extended by several independent works:

\begin{itemize}[leftmargin=2em]
  \item \textbf{Inference-Time Intervention (ITI)} \cite{li2023iti}:
    Intervenes on the activations of specific attention heads at inference time,
    significantly improving model truthfulness.
    ITI's findings are consistent with RepE: truthfulness has a localizable linear structure in model representations.

  \item \textbf{Activation Addition (ActAdd)} \cite{turner2023actadd}:
    Controls the generation direction by adding steering vectors during the model's forward pass.
    Technically highly similar to RepE's method, though ActAdd's steering vectors
    are typically obtained through simpler means (e.g., direct prompt contrast).

  \item \textbf{Refusal as a Single Direction} \cite{arditi2024refusal}:
    Arditi et al.\ found that LLM refusal behavior
    \textbf{is mediated by a single direction in representation space}.
    By ablating this direction, the model's refusal capability can be completely shut off;
    by amplifying it, the model can be made to refuse any request.
    This finding explains from the representation level why safety training is so fragile ---
    if the entire refusal behavior depends on only one linear direction,
    then any attack capable of interfering with this direction can breach safety alignment.

  \item \textbf{Circuit Breakers} \cite{zou2024circuitbreakers}:
    Zou et al.\ applied RepE's ideas to safety defense,
    training models to ``circuit break'' when unsafe representation directions are detected ---
    that is, mapping hidden states to a safe region to prevent the generation of unsafe content.
    Unlike traditional refusal training, Circuit Breakers operate at the representation level rather than the behavioral level,
    providing stronger robustness against unseen attacks.
\end{itemize}

\begin{solutionbox}
Representation engineering brings a fundamental shift in perspective to safety alignment:

\textbf{From ``training what the model says'' to ``understanding how the model thinks.''}

Traditional methods (CAI, Safe RLHF, red teaming) all operate at the behavioral level ---
modifying training data or objectives to change the model's input-output mapping.
RepE and its related works demonstrate that safety-related concepts (such as honesty, harmlessness, refusal)
possess \textbf{interpretable linear structure} in the model's representation space,
which can be directly read and manipulated.
This not only provides new control mechanisms for safety (e.g., Circuit Breakers)
but also offers theoretical tools for understanding why safety training works (or why it fails).
\end{solutionbox}

\evolution{Behavioral-level safety training (CAI, RLHF)}{Representation-level understanding and control (RepE, Circuit Breakers)}{Deeper safety assurance mechanisms are needed}

% ============================================================
\subsection{2025--2026: Safety Challenges of Reasoning Models}
\label{sec:t3-reasoning-safety}
% ============================================================

The proliferation of thinking models (\crossref{4}{sec:t4-thinking-models}) has introduced entirely new attack surfaces and challenge dimensions for safety alignment.
Once models possess autonomous reasoning and planning capabilities, safety concerns escalate from ``models passively complying with harmful instructions''
to ``models actively reasoning about how to circumvent safety restrictions.''

\subsubsection{Reasoning Models as Autonomous Jailbreak Agents}

Hagendorff et al.\ \cite{hagendorff2026jailbreakagents} published a study in Nature Communications revealing an alarming phenomenon:
\textbf{large reasoning models can autonomously act as jailbreak agents},
spontaneously constructing attack strategies to breach other models' safety defenses without the need for carefully crafted human prompts.
In experiments, the success rate of reasoning models acting as attack agents reached as high as \textbf{97\%},
far exceeding traditional manually designed jailbreak methods.

The far-reaching implication of this finding is that
\textbf{the safety alignment landscape is shifting from ``human attacker vs.\ AI defender''
to ``AI attacker vs.\ AI defender''}.
When the attack side is also powered by powerful reasoning models,
the level of robustness required for defense methods rises dramatically.

\subsubsection{Role-Playing and Multilingual Safety Vulnerabilities}

Tang et al.\ \cite{tang2025rolebreak} proposed RoleBreak, revealing safety risks of reasoning models from another angle:
through \textbf{character hallucination} in role-playing scenarios,
attackers can induce models to generate harmful content during the ``acting in character'' reasoning process,
bypassing standard safety training.
This class of attacks exploits a structural weakness of reasoning models:
the conflict between maintaining character consistency during long-chain reasoning and adhering to safety restrictions.

\subsubsection{Lifelong Safety Alignment}

Wang et al.\ \cite{wang2025lifelongsafety} proposed the \textbf{Lifelong Safety Alignment} framework,
addressing a problem that is particularly acute in the era of reasoning models:
does the safety alignment previously learned by models gradually erode
during continual fine-tuning and RL training?
This work proposed methods for maintaining safety alignment throughout a model's entire lifecycle,
including safety-aware continual learning and alignment anchoring mechanisms.

\subsubsection{2026: Comprehensive Deepening of Safety Alignment}

In early 2026, safety alignment research advanced simultaneously across the four dimensions of \textbf{attack}, \textbf{defense}, \textbf{theory}, and \textbf{evaluation},
exhibiting unprecedented systematic deepening.

\paragraph{GRP-Obliteration: The ``double-edged sword'' nature of alignment techniques.}
\landmark{russinovich2026grp}
Russinovich et al.\ from Microsoft Research revealed a disturbing finding:
\textbf{GRPO --- the very RL technique used for safety alignment ---
can be reverse-engineered to destroy a model's safety alignment using only a single unlabeled harmful prompt}.
Evaluation across 15 models (7B--20B) showed that
GRP-Obliteration's unalignment effect is stronger than all existing methods
while preserving the model's general capabilities.
Even more concerning, the safety vulnerability caused by a single harmful prompt
\textbf{spreads across all safety categories} ---
the model's entire safety guardrail is systematically dismantled.
This result fundamentally challenges the robustness assumption of RLHF/GRPO-based safety training.

\paragraph{Systematic evaluation and defense of reasoning model safety.}
Li et al.\ \cite{li2026whatmatters} conducted the largest-scale safety evaluation to date
across 32 models (3B--235B) (4.6M API calls, 56 jailbreak techniques),
revealing key factors influencing reasoning model safety:
\textbf{chain-of-thought response-prefix attacks}
increase the average attack success rate by 3.34$\times$,
with one model's vulnerability surging from 0.6\% to 96.3\%.
The study also found that reasoning-based self-reflection can improve safety,
but post-training distillation may \emph{silently} erode this capability.

On the defense side, SafeThink \cite{ghosal2026safethink}
proposed an \textbf{inference-time safety recovery} mechanism:
monitoring reasoning traces via a safety reward model
and injecting corrective prompts (``Wait, think safely'') when a safety threshold breach is detected.
The key insight is that \textbf{intervention in only the first 1--3 reasoning steps suffices}
to redirect the subsequent full generation onto a safe trajectory,
reducing attack success rates by 30--60\% without sacrificing reasoning quality.

\paragraph{Theoretical foundations and mitigation of the Alignment Tax.}
Alignment tax --- the capability loss caused by safety training ---
has evolved in 2026 from scattered empirical observations into a \textbf{systematic theoretical framework}.
Young \cite{young2026alignmenttax} provided a formal definition from a geometric perspective:
alignment tax equals the \textbf{squared projection} of the safety direction onto the capability subspace,
and derived a tight Pareto frontier for the safety-capability trade-off.
More importantly, the work proved that alignment tax can be decomposed into
an \emph{irreducible structural component} (determined by subspace angles) and
a \emph{packing residual} that decreases with model dimensionality ---
implying that larger models inherently possess lower alignment tax due to their geometric structure.

On the practical mitigation front,
Xie et al.\ \cite{xie2026dgr} identified the
\textbf{distributional mismatch} between safety data and the target model's distribution
as the root cause of reasoning capability degradation,
and proposed Distribution-Grounded Refinement (DGR),
which transforms the safety dataset to the model's native distribution before training,
improving reasoning accuracy by 30.2\%.
A surprising finding was that
\textbf{as few as 10 samples suffice to activate effective refusal behavior},
suggesting that safety functionality is more akin to knowledge \emph{activation} than extensive retraining.
Wang et al.\ \cite{wang2026ascl} reframed safety alignment as an
\textbf{adaptive retrieval} problem:
the model autonomously decides when to consult safety rules,
using Inverse Frequency Policy Optimization to counteract the capability degradation caused by excessive consultation.

\paragraph{Mechanisms and defenses against fine-tuning attacks.}
Zhang et al.\ \cite{zhang2026spf} revealed the \textbf{structural mechanisms}
underlying why fine-tuning attacks \cite{qi2023finetuning} work:
(1) safety gradients occupy a low-rank subspace, while utility gradients span higher dimensions;
(2) the two subspaces frequently exhibit directional conflicts;
(3) dominant safety directions can be estimated from very few samples.
Based on this, Safety-Preserving Fine-Tuning (SPF)
removes gradient update components that conflict with the safety subspace,
and provides convergence proofs with bounded safety drift.
This work forms a \textbf{theoretical corroboration} with Young's \cite{young2026alignmenttax} geometric analysis:
the low-rank structure of safety is both the source of vulnerability (easily destroyed by fine-tuning)
and the lever for protection (can be precisely estimated and projection-protected).

\paragraph{Innovation in safety evaluation methodology.}
Chen et al.\ \cite{chen2026expectedharm} challenged the current severity-based safety evaluation paradigm,
proposing the \textbf{Expected Harm} framework:
weighting severity by \emph{execution likelihood} (the conditional probability of threat realization).
This work discovered the \textbf{Inverse Risk Calibration} phenomenon:
models strongly refuse threats with low execution probability
yet remain vulnerable to attacks with high execution probability.
Exploiting this calibration bias doubles the jailbreak success rate.

\paragraph{New attack surfaces in multimodal safety.}
SIVA \cite{rashid2026siva} discovered that VLM safety alignment is conducted only on \textbf{complete images},
and completely fails when harmful semantics are \emph{distributed across multiple image fragments} ---
revealing a fundamental compositionality deficiency in visual safety training.
AM3Safety \cite{zhu2026am3safety} was the first to systematically address
the joint \textbf{multimodal + multi-turn} safety problem:
combining cold-start refusal with GRPO fine-tuning
and using a turn-aware dual-objective reward,
it reduced attack success rate by 10\% while improving the harmless (+8\%) and helpful (+13\%) dimensions
(\crossref{5}{sec:t5}).

\evolution{Behavioral-level safety training (RLHF/CAI) $\to$ representation-level understanding (RepE, LAT)}{Geometric theorization (Pareto frontier of alignment tax) + RL double-edged sword awareness (GRP-Obliteration)}{Safety alignment must transition from empirical tuning to theoretically grounded engineering}

% ============================================================
\subsection{Open Problems}
\label{sec:t3-open}
% ============================================================

Despite the tremendous progress in safety alignment over the past three years ---
from CAI's principle-driven methods to RepE's representation engineering, and further to the frontiers of reasoning model safety ---
the field still faces several fundamental open problems.

\begin{openproblembox}
\textbf{Open Problem 1: Safety--Capability Parity.}
The mismatched generalization framework of Wei et al.\ \cite{wei2023jailbroken} shows that
every time a model acquires a new capability (e.g., supporting a new language, understanding a new modality),
a new attack surface may open up.
How can safety scale in tandem with capability?
One possible direction is to deeply integrate safety training into the pre-training stage
(rather than treating it as an add-on during post-training),
but this would significantly increase pre-training costs and may affect the model's general capabilities.
\end{openproblembox}

\begin{openproblembox}
\textbf{Open Problem 2: Scalable Oversight.}
As model capabilities surpass those of humans,
traditional RLHF and CAI methods (both relying on human or existing AI judgments as supervisory signals)
will face fundamental limitations.
The theoretical roots of this problem can be traced back to
\textbf{iterated amplification} proposed by Christiano et al.\ \cite{christiano2018oversight} ---
decomposing complex tasks into subtasks that can be judged by weaker supervisors to indirectly supervise stronger models ---
as well as \textbf{AI safety via debate} proposed by Irving et al.\ \cite{irving2018debate} ---
having two AI systems debate each other, with humans needing only to judge the persuasiveness of the final arguments.
Weak-to-strong generalization \cite{burns2023weaktostrong} provides preliminary empirical evidence
suggesting that weaker supervisors can guide stronger models to some extent,
but the reliability of this approach in safety-critical scenarios has yet to be verified.
\end{openproblembox}

\begin{openproblembox}
\textbf{Open Problem 3: Deceptive Alignment.}
Sleeper Agents \cite{hubinger2024sleeper} demonstrated that current safety training
cannot reliably detect and eliminate hidden unsafe behavior in models.
If future AI systems spontaneously learn strategic deception during training
(strategically deceiving evaluators to appear aligned),
all existing behavioral-level evaluation methods will fail.
Representation-level methods such as RepE offer a theoretical possibility for detecting such behavior,
but current techniques are far from mature enough to provide reliable guarantees.
\end{openproblembox}

\begin{openproblembox}
\textbf{Open Problem 4: Robustness of Representation-Level Safety.}
RepE and Circuit Breakers have demonstrated the potential of controlling safety behavior at the representation level,
but the findings of Arditi et al.\ \cite{arditi2024refusal} also imply a risk:
if refusal is encoded by only a single linear direction,
can representation-level defenses be equally easily bypassed by representation-level attacks?
Will adversarial representation manipulation
become the primary direction of next-generation jailbreaking?
\end{openproblembox}

\begin{openproblembox}
\textbf{Open Problem 5: Reasoning Chain Safety.}
The hidden CoT of reasoning models introduces an entirely new dimension of safety challenges:
how can we ensure that reasoning chains themselves do not contain harmful reasoning
(e.g., ``Let me think about how to bypass safety restrictions...''),
while not excessively constraining the model's reasoning capabilities?
Monitoring and auditing hidden CoT is technically feasible,
but imposing content restrictions on reasoning chains may impair reasoning quality.
The more fundamental question is: when reasoning models themselves become jailbreak agents
\cite{hagendorff2026jailbreakagents},
is traditional input/output-level safety training fundamentally insufficient?
Is it necessary to enforce safety constraints at the reasoning process level (rather than only at the output level)?
\end{openproblembox}

\begin{openproblembox}
\textbf{Open Problem 6: Multimodal Safety.}
When safety research extends from text-only to multimodal settings,
the attack surface expands dramatically --- images, audio, and video can all serve as attack vectors.
Visual adversarial examples can embed instructions in images that are invisible to humans but visible to models.
A systematic solution for multimodal safety remains an open problem
(\crossref{5}{sec:t5}).
\end{openproblembox}

\subsection{Chapter Summary}

\begin{figure}[!ht]
\centering
\begin{tikzpicture}[
  node distance=0.6cm and 1.2cm,
  every node/.style={font=\small},
  milestone/.style={rectangle, rounded corners, draw=thread3, fill=thread3!10,
    text width=4.2cm, minimum height=0.8cm, align=center, font=\small\bfseries},
  branch/.style={rectangle, rounded corners, draw=gray!60, fill=gray!5,
    text width=3.2cm, minimum height=0.7cm, align=center, font=\small},
  arrow/.style={-{Stealth[length=2.5mm]}, thick, thread3!70},
  brancharrow/.style={-{Stealth[length=2mm]}, gray!60, thick},
]
% Main timeline
\node[milestone] (hh) {HH-RLHF (2022)\\Helpfulness--Harmlessness Tension};
\node[milestone, below=1.2cm of hh] (cai) {Constitutional AI (2022)\\AI Self-Alignment};
\node[milestone, below=1.2cm of cai] (redteam) {Red Teaming (2022)\\LLM-based Adversarial Testing};

% Branches
\node[branch, below left=1.2cm and 0.5cm of redteam] (brA) {Branch A: Training\\Safe RLHF, Lagrangian};
\node[branch, below right=1.2cm and 0.5cm of redteam] (brB) {Branch B: Attack Analysis\\GCG, Jailbreak Taxonomy};

% Convergence
\node[milestone, below=2.8cm of redteam] (repe) {RepE (2023)\\Representation Engineering};
\node[milestone, below=1.2cm of repe] (agent) {Reasoning Model Safety\\(2025--2026)};
\node[branch, below=1.2cm of agent] (conv) {Paradigm Shift: From Behavioral\\to Representation + Agent Safety};

% Arrows
\draw[arrow] (hh) -- (cai);
\draw[arrow] (cai) -- (redteam);
\draw[brancharrow] (redteam) -- (brA);
\draw[brancharrow] (redteam) -- (brB);
\draw[arrow] (redteam) -- (repe);
\draw[arrow] (repe) -- (agent);
\draw[brancharrow] (brA) -- (repe);
\draw[brancharrow] (brB) -- (repe);
\draw[arrow] (agent) -- (conv);
\end{tikzpicture}
\caption{Thread 3 evolutionary trajectory: from helpfulness--harmlessness tension to representation engineering and agent safety.
Red nodes are landmark works; gray nodes are important branches.}
\label{fig:safety-evolution}
\end{figure}

\begin{table}[!ht]
\centering
\caption{Comparison of core methods in Thread 3. Category distinguishes training-stage methods, attack methods, and analysis/control methods.}
\label{tab:safety-comparison}
\small
\begin{tabularx}{\textwidth}{l c c >{\raggedright\arraybackslash}X >{\raggedright\arraybackslash}X}
\toprule
\textbf{Method} & \textbf{Category} & \textbf{Human Labels?} & \textbf{Key Mechanism} & \textbf{Limitation} \\
\midrule
HH-RLHF & Training & Yes (pairwise) & Single reward encoding helpfulness + harmlessness & Evasiveness from objective conflict \\
CAI & Training & No (AI) & Constitution principles + RLAIF self-critique & Bounded by model's own judgment \\
Safe RLHF & Training & Yes (decoupled) & CMDP: maximize $R$ s.t.\ $C \leq d$; Lagrangian $\lambda$ & Still requires human annotation \\
Red Teaming & Evaluation & Minimal & LM generates attacks; RL-optimized for success rate & Coverage limited to known patterns \\
GCG & Attack & None & Gradient-based adversarial suffix optimization & White-box; but cross-model transfer \\
RepE & Control & No & Linear steering vectors in representation space & Safety encoded as fragile single direction \\
\bottomrule
\end{tabularx}
\end{table}

The evolution of Thread~3 exhibits a clear main thread:
from HH-RLHF exposing the helpfulness--harm\-less\-ness tension,
through CAI and Safe RLHF providing systematic training solutions,
to the jailbreaking arms race continuously challenging the robustness of these solutions.
RepE initiated a paradigm shift from the behavioral level to the representation level,
while the proliferation of reasoning models in 2025--2026 elevated safety challenges to an entirely new dimension ---
when models themselves become attack agents, the entire adversarial landscape undergoes a fundamental transformation.
This thread has close cross-cutting connections with Thread~2 (preference optimization provides the foundational framework, \crossref{2}{sec:t2}),
Thread~4 (reasoning capabilities affect the reliability of models following safety reasoning, \crossref{4}{sec:t4}),
and Thread~5 (multimodal extensions introduce new safety challenges, \crossref{5}{sec:t5}).
Safety alignment is not a static problem that can be ``solved,''
but rather a dynamic research direction that continually evolves alongside model capabilities.

\paragraph{Preview of the next chapter.}
Safety research has revealed a key fact: the stronger a model's reasoning capabilities, the more complex its safety challenges become.
Thread~4 will trace the evolution of reasoning capabilities ---
from chain-of-thought prompting to thinking models becoming the default paradigm,
showing how post-training internalizes reasoning capabilities into model parameters through RL and process supervision.
